\chapter{Parents' Self-Empathy and Self-Regulation}

\begin{quote}
\itshape ``If you are not able to be empathetic towards yourself,
you will not be able to be empathetic towards anyone else.''\\
\upshape --- Marshall B. Rosenberg
\end{quote}

\section{The Foundation of Self-Regulation}

Before we can effectively connect with our children's emotions,
we must first be able to recognize and regulate our own.
\citet{siegel2003} emphasize that a parent's ability to remain
calm and integrated in the face of a child's distress is
fundamental to the child's own emotional development.
When we are triggered, our nervous system enters a state of
fight, flight, or freeze, making it physiologically difficult
to access the prefrontal cortex---the area of the brain
responsible for empathy, logic, and connection.

Our children rely on us to lend them our calm nervous
system. When a child is upset, they look to the parent to
coregulate. If the parent responds with anxiety, anger, or
frustration, the child's distress is often amplified.
Conversely, when a parent can ground themselves and remain
centered, they provide a safe container for the child's
big emotions. This process of coregulation helps the child
learn, over time, how to self-regulate \citep{porges2011}.

Recognizing our own triggers is the first step. Triggers are
often rooted in our own unmet needs or past experiences.
When a child's behavior elicits a disproportionate reaction
in us, it is usually a signal that an old wound has been
touched. By bringing awareness to these moments, we can begin
to pause and choose our response, rather than reacting
automatically.

\begin{keyinsight}
A parent's self-regulation is the prerequisite for coregulating
with a child. When parents manage their own nervous system
arousal, they model emotional regulation and provide the safety
necessary for children to process their own emotions.
\end{keyinsight}

\section{Cultivating Self-Empathy}

Self-empathy is the practice of turning the compassionate
attention we usually reserve for others inward toward ourselves.
In the framework of Nonviolent Communication, \citet{rosenberg2003}
teaches that self-empathy involves observing our own feelings and
needs without judgment, blame, or evaluation.

When we make a mistake or react in a way we regret, our habitual
response is often self-criticism. We might think, ``I shouldn't
have yelled,'' or ``I am a terrible parent.'' These thoughts
create further disconnection from ourselves and, consequently,
from our children. Self-empathy invites us to replace
self-judgment with curiosity. We might ask ourselves, ``What am
I feeling right now?'' and ``What need of mine is not being met?''

Perhaps the yelling was fueled by a feeling of exhaustion and a
need for rest. Or perhaps it stemmed from a feeling of overwhelm
and a need for support. By identifying the underlying feelings
and needs, we shift our focus from what is ``wrong'' with us to
what is alive in us. This shift softens our internal landscape,
making it easier to forgive ourselves and repair the connection
with our child.

\begin{keyinsight}
Self-empathy involves observing our own feelings and needs without
judgment or blame. By shifting from self-criticism to curious
exploration of our unmet needs, we cultivate the inner resilience
required for empathetic parenting.
\end{keyinsight}

\section{Practice Exercise}

\begin{practiceexercise}
\textbf{Multiple Choice Questions}
\begin{enumerate}
  \item When a parent is triggered and enters a fight, flight, or freeze response, which part of the brain becomes difficult to access?
  \begin{enumerate}
    \item The amygdala.
    \item The brainstem.
    \item The prefrontal cortex.
    \item The cerebellum.
  \end{enumerate}
  \vspace{\fill}
  \item According to Nonviolent Communication, what is the alternative to self-criticism when a parent makes a mistake?
  \begin{enumerate}
    \item Ignoring the mistake to maintain confidence.
    \item Blaming the child for triggering the reaction.
    \item Asking others for validation.
    \item Observing one's own feelings and unmet needs with curiosity.
  \end{enumerate}
  \vspace{\fill}
  \item What is the process called when a calm parent helps an upset child manage their emotional state?
  \begin{enumerate}
    \item Discipline.
    \item Coregulation.
    \item Distraction.
    \item Enmeshment.
  \end{enumerate}
\end{enumerate}
\vspace{\fill}
\newpage

\textbf{Explanations for Questions 1--3}

\textbf{1. (c) The prefrontal cortex.} When the nervous system is highly aroused by a trigger, the brain prioritizes survival mechanisms, often temporarily impairing access to the prefrontal cortex, which governs empathy, logic, and connection.

\textbf{2. (d) Observing one's own feelings and unmet needs with curiosity.} Self-empathy shifts the focus away from judgment and blame, encouraging a compassionate exploration of the underlying emotions and needs that drove the behavior.

\textbf{3. (b) Coregulation.} Coregulation occurs when a regulated adult provides a safe, calming presence that helps soothe an upset child's nervous system, teaching them how to eventually self-regulate.
\vspace{\fill}
\newpage

\begin{enumerate}
  \setcounter{enumi}{3}
  \item A parent's disproportionate reaction to a child's behavior is often a signal of:
  \begin{enumerate}
    \item The child's intentional defiance.
    \item A lack of discipline in the household.
    \item An old wound or unmet need being touched in the parent.
    \item The child's inability to understand rules.
  \end{enumerate}
  \vspace{\fill}
  \item Which of the following is an example of a self-empathetic thought?
  \begin{enumerate}
    \item "I shouldn't be so tired all the time."
    \item "I am feeling overwhelmed because I have a need for support right now."
    \item "I need to try harder to be a perfect parent."
    \item "If my child behaved better, I wouldn't be so angry."
  \end{enumerate}
  \vspace{\fill}
  \item Why is parental self-regulation considered a prerequisite for helping a child process emotions?
  \begin{enumerate}
    \item Because children will only listen to parents who never show emotion.
    \item Because an anxious or angry parent can amplify the child's distress, whereas a calm parent provides a safe container.
    \item Because parents need to prove they are in charge.
    \item Because children are incapable of having emotions without adult instruction.
  \end{enumerate}
\end{enumerate}
\vspace{\fill}
\newpage

\textbf{Explanations for Questions 4--6}

\textbf{4. (c) An old wound or unmet need being touched in the parent.} Triggers often reflect our own past experiences or current unmet needs rather than solely being a response to the child's immediate action.

\textbf{5. (b) "I am feeling overwhelmed because I have a need for support right now."} This statement neutrally observes a feeling and connects it to a specific, underlying human need without resorting to judgment or blame.

\textbf{6. (b) Because an anxious or angry parent can amplify the child's distress, whereas a calm parent provides a safe container.} Children rely on their parents' nervous systems to help them regulate; a dysregulated parent cannot effectively soothe a dysregulated child.
\vspace{\fill}
\newpage

\begin{enumerate}
  \setcounter{enumi}{6}
  \item In the context of self-regulation, what does "lending our calm nervous system" mean?
  \begin{enumerate}
    \item Ignoring the child's feelings so they learn to cope alone.
    \item Pretending to be happy even when feeling furious.
    \item Maintaining our own groundedness to provide a stabilizing presence for the child.
    \item Telling the child to calm down immediately.
  \end{enumerate}
  \vspace{\fill}
  \item How does self-empathy help repair connection with a child after a rupture?
  \begin{enumerate}
    \item It makes the parent realize the child was right.
    \item It forces the parent to apologize profusely to avoid guilt.
    \item It softens the parent's internal landscape, reducing defensiveness and making genuine repair easier.
    \item It convinces the child that the parent's behavior was justified.
  \end{enumerate}
  \vspace{\fill}
  \item What is the ultimate goal of recognizing our own triggers?
  \begin{enumerate}
    \item To eliminate all negative emotions from our lives.
    \item To be able to pause and choose our response rather than reacting automatically.
    \item To train our children to avoid those triggers.
    \item To justify our reactions to stressful situations.
  \end{enumerate}
\end{enumerate}
\vspace{\fill}
\newpage

\textbf{Explanations for Questions 7--9}

\textbf{7. (c) Maintaining our own groundedness to provide a stabilizing presence for the child.} When we stay regulated, our calm state can literally help soothe the child's aroused nervous system through the process of coregulation.

\textbf{8. (c) It softens the parent's internal landscape, reducing defensiveness and making genuine repair easier.} By understanding our own needs with compassion, we release self-blame, which opens the door to authentic connection and apology with our child.

\textbf{9. (b) To be able to pause and choose our response rather than reacting automatically.} Awareness of our triggers creates the necessary space between stimulus and response, allowing us to parent aligned with our values.
\vspace{\fill}
\newpage
\end{practiceexercise}

\section{Reflection Questions}

\begin{reflection}
\small
\begin{enumerate}[itemsep=0pt, parsep=0pt, topsep=0pt]
  \item When you feel triggered by your child's behavior, where do you typically feel it in your body first?
  \item What specific behaviors from your child tend to evoke the strongest emotional reactions in you?
  \item Looking back at your own childhood, how did your caregivers handle their own stress or anger?
  \item In moments of parenting stress, what unmet needs (e.g., rest, support, understanding) are usually present for you?
  \item When you make a parenting mistake, what does your inner monologue sound like (e.g., critical, compassionate, dismissive)?
  \item What is one small, realistic action you can take to pause and ground yourself before reacting to a trigger?
  \item How might approaching yourself with the same empathy you wish to give your child change the atmosphere in your home?
  \item Can you recall a time when you successfully coregulated with your upset child? What helped you stay calm?
  \item Who or what provides you with the emotional support you need to maintain your own self-regulation?
\end{enumerate}
\end{reflection}